\scanpage{572}
\mspage{481}
ordinaux $\alpha$ \deleted{\unsure{$\in$}} qui sont $\in\unsure{\mathcal U|U}$. Donc $I=I_{\alpha_0}$, où
$\alpha_0$ est le plus petit ordinal qui n'appartient
pas à $\mathcal U$, et $I$ est grand devant tt
\inserted{petit} cardinal \deleted{\unsure{de}} \unsure{i.e.|de} devant tt cardinal $\in\mathcal U$.
% slanted
\margin{introduire $I=I_{\mathcal U}$ avant l'énoncé de la proposition !}
La condition envisagée est la suivante :

Il existe une suite transfinie
\[
\underline x=(x_\alpha)_{\alpha\in I}
\]
d'éléments de $\operatorname{Ob}M$, telle que l'on ait
\[
\operatorname{Hom}(x_{\deleted{\unsure{\alpha}}\beta},x_{\deleted{\unsure{\beta}}\alpha})=\varnothing\qquad\text{si }\beta<\deleted{\unsure{\beta}}\alpha.
\]

De façon précise, pour une telle suite
transfinie, soit $\underline S\deleted{\unsure{\subset}}\unsure{=}\,{}_{\underline x}M$ la sous-catégorie
fermée de \deleted{\unsure{$\underline S$}} $M$ engendrée par la
famille des $x_\alpha$, \inserted{i.e telle que $\operatorname{Ob}\underline S=\bigcup_{\alpha\in I}\operatorname{Ob}{}_{x_\alpha}M=\bigcup_{\alpha\in I}x_\alpha\backslash\operatorname{Ob}M$.}
Alors $ {}_{\underline x}M$ est une
sous-catégorie fermée non accessible dans $M$,
et toute \deleted{\unsure{telle}} sous-catégorie fermée non
accessible dans $M$ s'obtient de cette
façon.

\uline{Démonstration} Considérons une suite
$\underline x=(x_\alpha)_{\alpha\in I}$ comme dessus, \deleted{\unsure{et la}}
\deleted{En vertu de 4.16.7.1,} si $J$ est une
petite partie de $I$, \deleted{et si \unsure{$\alpha$}} il existe
un $\alpha\in I$ qui majore les $\beta\in J$,
donc $\operatorname{Hom}(x_\beta,x_\alpha)=\varnothing$ pour $\beta\in J$,
% archive sheet number: 571
